% =============================================================================
%  Experimental Setup
% =============================================================================
% Requires in the parent document:
% \usepackage{booktabs}

\section{Experimental Setup}
\label{sec:experimental_setup}

All solver experiments were run on the same local machine. We did not use a
cloud instance, a separate server, or GPU acceleration for the solver runs. This
keeps the comparison grounded in one environment: the same CPU, memory, operating
system, and benchmark runner for every solver.

\subsection{Machine Configuration}

The machine used for the experiments is listed in
Table~\ref{tab:experimental_machine_setup}.

\begin{table}[htbp]
  \centering
  \caption{Hardware and operating system used for the experiments}
  \label{tab:experimental_machine_setup}
  \begin{tabular}{ll}
    \toprule
    Component & Specification \\
    \midrule
    Machine & Acer Predator PH315-55 \\
    System type & x64-based PC \\
    Processor & 12th Gen Intel(R) Core(TM) i7-12700H \\
    CPU cores / logical processors & 14 cores / 20 logical processors \\
    Maximum reported CPU clock & 2700 MHz \\
    Installed memory & 16 GB physical memory \\
    Operating system & Microsoft Windows 11 IoT Enterprise LTSC \\
    OS version & 10.0.26100, build 26100, 64-bit \\
    Graphics adapters & Intel(R) Iris(R) Xe Graphics; NVIDIA GeForce RTX 3060 Laptop GPU \\
    \bottomrule
  \end{tabular}
\end{table}

The solvers themselves run on the CPU. The machine has both integrated and
dedicated graphics adapters, but the benchmark code does not use them for
solver computation.

\subsection{Software Environment}

The project is a Python project. Its configuration requires
Python \(3.13 \leq \text{version} < 3.14\). The runtime dependencies are NumPy,
Pygame, Z3 Solver, Matplotlib, Seaborn, and Pandas. Pytest is included for
development and testing.

The benchmark runner lives in \texttt{src/benchmark/benchmark.py}. The plots are
generated by \texttt{src/benchmark/visualize.py}. For each run, the benchmark
code records the solver name, puzzle size, input file name, success flag, status
message, elapsed time in milliseconds, peak memory in KB, inference count, node
expansion count, backtrack count, and completion ratio.

\subsection{Benchmark Data}

The benchmark puzzles are stored in \texttt{src/benchmark/input}. Their expected
solutions are stored in \texttt{src/benchmark/expected}. There are 24 puzzles in
total, covering sizes from \(4 \times 4\) to \(9 \times 9\). Each size has four
puzzles. By difficulty, the set contains six easy puzzles, twelve medium
puzzles, and six hard puzzles.

\begin{table}[htbp]
  \centering
  \caption{Benchmark puzzle distribution}
  \label{tab:benchmark_puzzle_distribution}
  \begin{tabular}{lrrrrrr}
    \toprule
    Puzzle size & \(4\times4\) & \(5\times5\) & \(6\times6\) & \(7\times7\) & \(8\times8\) & \(9\times9\) \\
    \midrule
    Number of puzzles & 4 & 4 & 4 & 4 & 4 & 4 \\
    \bottomrule
  \end{tabular}
\end{table}

\subsection{Solvers Tested}

The benchmark runner can run the following solver keys:
\texttt{forward\_chaining},
\texttt{forward\_then\_backward\_chaining},
\texttt{backtracking\_forward\_chaining},
\texttt{backward\_chaining},
\texttt{brute\_force},
\texttt{astar\_h1},
\texttt{astar\_h2},
\texttt{astar\_h3}, and
\texttt{astar\_h4}.

The A* variants use four heuristics: empty cells, domain sum, min-conflicts, and
AC-3 pruned domain sum. All solvers go through the same benchmark path. The
runner parses the puzzle, runs the selected solver, checks the produced grid
against the expected solution, and records the same measurement fields.

\subsection{Measurement Procedure}

One benchmark case means one solver running on one puzzle file. Runtime is stored
in milliseconds. Peak memory is stored in KB. A* solvers use node expansions as
their operation count, while the chaining solvers use inference count. The
runner also records completion ratio, which makes partial outputs easier to
compare in the tables and plots.

Plots are generated only after the benchmark rows have already been collected.
They are a presentation step, not part of the solver execution.

\subsection{Statistics Collection Method}

The benchmark runner collects statistics directly from each solver run. For a
single case, it parses the puzzle, creates the selected solver, runs that solver
on a copy of the puzzle, and reads the \texttt{Stats} object returned at the end.
Using a copy matters because it keeps one solver from changing the puzzle object
seen by another solver.

Runtime is measured with Python's \texttt{time.perf\_counter()} and then
converted to milliseconds. Memory is measured with Python's \texttt{tracemalloc}
module. Each solver records the peak traced memory during its solve step and
converts the value from bytes to KB. This is Python-traced memory for the solver
run, not the full memory footprint of the operating system process.

The operation fields depend on the solver type. A* solvers fill
\texttt{node\_expansions}, since their work is based on expanding search states.
Chaining solvers fill \texttt{inference\_count}, since their work is based on
applying rules. The backtracking-forward-chaining solver also records
\texttt{backtracks}. If a metric does not apply to a solver, the value is left
as zero.

After a solver finishes, the runner compares the produced grid with the expected
solution file of the same name. It records whether the solution is complete,
whether the grid matches, and the \texttt{completion\_ratio}. The completion
ratio is the fraction of initially empty cells that the solver filled.

The runner then writes one row for that puzzle case. The row includes the test
name, solver name, puzzle size, input file, success flag, status message,
runtime, peak memory, operation counts, backtracks, and completion ratio. The
visualization script reads those rows later to create the runtime, memory,
operation-count, and completion-ratio plots.
